% Chapter 15 -- Registers, buses and the datapath
\chapter{Registers, buses and the datapath}\label{ch:datapath}

In ch14 the ALU answered questions. It could not \emph{remember}
anything --- feed it $3+5$ and it hands you 8, but the 8 evaporates the
moment you change the inputs. This chapter builds the machinery that
makes answers stick: a register file to hold intermediate values, the
bus wiring that lets blocks share one set of wires without fighting,
program memory that turns a list of instructions into data, and the
datapath itself --- the arrangement of all these pieces around the ALU.
The datapath is the figure this book has been building toward: once it
is on the page, ch16 only has to add the control unit that drives its
11-bit control word. By the end you will have watched a ten-byte
program execute, one control word at a time, on hardware you can
synthesise.

\section{The register file}\label{sec:regfile15}

\index{register file}\index{register!file}
A register file is a small array of identical registers with two kinds
of ports. The \emph{write port} accepts an address, a value, and a
clock edge --- the addressed cell latches the value then. The
\emph{read ports} accept only an address and answer combinationally,
all the time. Ours has eight 8-bit cells, one write port and two read
ports, because the ALU needs two operands simultaneously.

\begin{figure}[htbp]
\centering
\begin{tikzpicture}[font=\scriptsize, >=latex]
% 8 register cells in a column
\foreach \i in {0,...,7}{
  \pgfmathsetmacro{\y}{6.3-\i*0.75}
  \draw[thick, fill=graybg] (2.2,\y) rectangle (4.4,\y+0.6);
  \node at (3.3,\y+0.3) {r\i\ /8};
}
% write address decoder
\draw[thick, fill=amberbg] (0.2,2.4) rectangle (1.9,3.6);
\node[align=center] at (1.05,3.0) {decoder\\\texttt{i\_waddr}/3};
\foreach \i in {0,...,7}{
  \pgfmathsetmacro{\y}{6.6-\i*0.75}
  \draw[->, slate] (1.9,3.0) -- (2.2,\y);
}
\draw[->, thick] (0.2,4.3) -- (1.05,4.3) -- (1.05,3.6);
\node[left] at (0.2,4.3) {\texttt{i\_clk}};
\draw[->, thick] (0.2,5.1) -- (1.05,5.1) -- (1.05,4.4);
\node[left] at (0.2,5.1) {\texttt{i\_we}};
\draw[->, thick] (0.2,5.9) -- (2.2,5.9) node[pos=0.55, above]{\texttt{i\_waddr}/3};
\draw[->, thick] (0.2,0.4) -- (1.05,0.4) -- (1.05,1.8) -- (2.2,1.8);
\node[left] at (0.2,0.4) {\texttt{i\_wdata}/8};
% two async read ports: any cell -> either output
\foreach \i in {0,...,7}{
  \pgfmathsetmacro{\y}{6.6-\i*0.75}
  \draw[thick, accent] (4.4,\y) -- (5.6,\y);
}
\draw[thick, fill=goodbg] (5.6,4.0) rectangle (6.9,4.9);
\node[align=center] at (6.25,4.45) {mux A\\\texttt{i\_raddr\_a}};
\draw[->, thick, accent] (6.9,4.45) -- (7.9,4.45) node[right]{\texttt{o\_ra}/8};
\draw[thick, fill=goodbg] (5.6,2.4) rectangle (6.9,3.3);
\node[align=center] at (6.25,2.85) {mux B\\\texttt{i\_raddr\_b}};
\draw[->, thick, accent] (6.9,2.85) -- (7.9,2.85) node[right]{\texttt{o\_rb}/8};
\node[align=center, text=slate] at (3.3,-0.3) {write: clocked, one cell\\read: combinational, both ports at once};
\end{tikzpicture}

\caption{The 8$\times$8 register file. A decoder picks the cell the
write edge updates; two independent muxes let both ALU operands be
read at once. Reads are combinational --- no clock involved ---
because the ALU wants the operands \emph{before} the edge that will
write back its result. Regenerate: \ttfamily
figures/ch15/fig\_regfile.tex.}
\label{fig:regfile15}
\end{figure}

\begin{lstlisting}[language=Verilog, caption={\ttfamily code/ch15\_regfile8.v
--- one synchronous write port, two asynchronous read ports. The whole
module is twenty lines; the conventions carry the weight.},
label={lst:regfile8}]
module regfile8 (
  input  wire        i_clk,
  input  wire        i_we,
  input  wire [2:0]  i_waddr,
  input  wire [7:0]  i_wdata,
  input  wire [2:0]  i_raddr_a,
  input  wire [2:0]  i_raddr_b,
  output wire [7:0]  o_ra,
  output wire [7:0]  o_rb
);
  reg [7:0] mem [0:7];
  assign o_ra = mem[i_raddr_a];
  assign o_rb = mem[i_raddr_b];
  always @(posedge i_clk) begin
    if (i_we) mem[i_waddr] <= i_wdata;
  end
endmodule
\end{lstlisting}

Two subtleties to notice before you ever read a bigger one.
The writes are \emph{synchronous} --- the value lands on a clock edge,
and \texttt{i\_we} is what decides whether any cell changes. That
write-enable line is the single most important signal in the whole
machine: it is how the control unit says ``remember this'' versus
``compute but keep nothing'' --- the difference between SUB and CMP
from ch14. And the reads are \emph{asynchronous}: changing an address
changes the output immediately, which is why a combinational ALU can
consume them in the same cycle.

\begin{example}[The register file, exercised]
\texttt{ch15\_regfile8\_tb.v} writes all eight cells, reads them back
on both ports, and then checks the edge semantics directly --- driving
a write and verifying the read port still shows the \emph{old} value a
nanosecond before the edge:

\begin{lstlisting}[caption={Real run of \ttfamily ch15\_regfile8\_tb.v.}]
PASS: regfile8 12 checks clean
\end{lstlisting}

The ``pre-edge still old / post-edge new'' pair is the check that
matters most: it proves writes respect the clock. An accidental
combinational write-through would fail it instantly.
\end{example}

\section{The bus, and who may drive it}\label{sec:bus15}

\index{bus}\index{tri-state}
Wire the register file's output, the ALU's output, and a memory's
output all onto one 8-wire bundle and you have a \emph{bus}. But only
one of them may talk at a time, or the wires carry a fight between 0
and 1 that settles as neither --- and in real silicon, as heat.
Textbook CPUs solve this with \emph{tri-state} drivers: each talker
has an output-enable, and all but one output-enable are off, so the
silent ones are electrically absent, not just quiet.

On an FPGA there is no third state. Inside the fabric every bus is a
set of multiplexers: one line always drives each wire, and a select
input chooses which talker it is. The idea is the same --- one talker
at a time --- but the mechanism is a mux, not a driver that disappears.
That is why our datapath writes \texttt{i\_wb\_sel} and
\texttt{i\_dst\_sel} instead of output-enables.

\begin{figure}[htbp]
\centering
\begin{tikzpicture}[font=\scriptsize, >=latex]
% --- left: the tri-state idea ---
\node[font=\small\bfseries] at (1.9,6.6) {tri-state idea (breadboard)};
\draw[thick] (-0.4,1.2) -- (4.2,1.2);
\node[below] at (1.9,1.2) {shared bus /8};
\foreach \x/\nm in {0.2/{reg A},2.0/{reg B},3.6/{ALU}}{
  \draw[thick, fill=graybg] (\x-0.35,2.4) rectangle (\x+0.35,3.4);
  \node at (\x,2.9) {\nm};
  \draw[thick, accent] (\x,2.4) -- (\x,1.75);
  \draw[thick, accent] (\x,1.2) -- (\x,1.75);
  \draw[thick, fill=amberbg] (\x-0.28,1.75) rectangle (\x+0.28,2.05);
  \node[font=\tiny] at (\x,1.9) {OE};
}
\node[align=center, text=slate] at (1.9,0.4) {each driver \emph{electrically disconnects}\\
when off --- two drivers on = smoke};
% --- right: the mux reality ---
\node[font=\small\bfseries] at (8.2,6.6) {mux reality (FPGA fabric)};
\draw[thick, fill=graybg] (6.6,2.4) rectangle (7.3,3.4);
\node at (6.95,2.9) {A};
\draw[thick, fill=graybg] (7.9,2.4) rectangle (8.6,3.4);
\node at (8.25,2.9) {B};
\draw[thick, fill=graybg] (9.2,2.4) rectangle (9.9,3.4);
\node[font=\tiny] at (9.55,2.9) {ALU};
\draw[thick, fill=accentsoft] (7.4,1.0) rectangle (8.8,1.9);
\node[align=center] at (8.1,1.45) {3:1 mux\\\texttt{sel}};
\draw[->, thick] (6.95,2.4) -- (6.95,1.45) -- (7.4,1.45);
\draw[->, thick] (8.25,2.4) -- (8.1,1.9);
\draw[->, thick] (9.55,2.4) -- (9.55,1.45) -- (8.8,1.45);
\draw[->, thick, accent] (8.1,1.0) -- (8.1,0.4);
\node[below, font=\small] at (8.1,0.4) {bus};
\node[align=center, text=slate] at (8.2,-0.5) {no third state exists: a select\\
line picks exactly one driver};
\end{tikzpicture}

\caption{Same idea, two mechanisms. Left: tri-state drivers on a
breadboard bus --- each component can electrically let go. Right: on
an FPGA the ``bus'' is a mux whose select line picks exactly one
driver. Both guarantee one talker; only the mux version synthesises.
Regenerate: \ttfamily figures/ch15/fig\_bus.tex.}
\label{fig:bus15}
\end{figure}

\begin{checkpoint}
Why doesn't the register file need an output-enable on its read
ports? Because read ports never drive the bus --- they are inputs to
muxes, not talkers on it. Only \emph{output} producers need steering.
\end{checkpoint}

\section{Memory: where the program lives}\label{sec:memory15}

\index{ROM}\index{RAM}\index{BRAM}
The datapath needs two kinds of memory with opposite personalities.
The \emph{program ROM} holds the instruction stream: read-only,
addressed by the program counter, loaded at build time. The \emph{data
RAM} holds scratch values: written and read by the running program.

\begin{lstlisting}[language=Verilog, caption={\ttfamily code/ch15\_rom.v
and \ttfamily ch15\_ram.v --- the two personalities side by side. The
ROM reads combinationally (fetch wants the instruction now); the RAM
registers its read output, the pattern that maps onto BRAM.},
label={lst:romram}]
// rom16: 16x8, combinational read -- the program is data in memory
reg [7:0] mem [0:15];
initial $readmemh("ch15_program.hex", mem);
assign o_data = mem[i_addr];

// ram16: 16x8, synchronous write AND read
always @(posedge i_clk) begin
  if (i_we) mem[i_addr] <= i_wdata;
  o_rdata <= mem[i_addr];
end
\end{lstlisting}

The RAM's extra \texttt{o\_rdata <=} line is not decoration. Block
RAMs inside the FPGA are synchronous memories: present an address and
the data arrives \emph{next} edge. Write the behavioural model with a
registered output and the synthesiser maps it to the RAM primitives;
write a combinational read and it may land on a pile of flops or
distributed LUT-RAM instead. The figures in the next section show what
actually happened to ours.

\begin{example}[Program ROM: the program is data]
\texttt{ch15\_mem\_tb.v} reads the demo program straight back out of
\texttt{rom16} --- the first byte at address 0 is \texttt{0x10},
our LDI opcode, and the byte at 1 is its immediate, \texttt{0x05}:

\begin{lstlisting}[caption={Real run of \ttfamily ch15\_mem\_tb.v ---
five ROM reads plus three RAM read/write/latency checks.}]
PASS: rom16 + ram16 8 checks clean
\end{lstlisting}

Sit with that for a second. The ROM does not know it holds a program.
Address in, byte out --- the \emph{meaning} of \texttt{0x10} is
created entirely by the pc stepping through addresses. This is the
von Neumann idea in miniature, and it is why ``a memory of
instructions read by an address counter'' is the whole trick of the
stored-program computer.
\end{example}

\begin{figure}[htbp]
\centering
\begin{tikzpicture}[font=\scriptsize, >=latex]
% program address space (left column)
\node[font=\small\bfseries] at (1.3,6.9) {program space};
\draw[thick] (0,0.4) rectangle (2.6,6.6);
\foreach \y/\lbl/\det in {6.0/{\texttt{0x00..0x09}}/{demo program\\(10 bytes)},3.6/{\texttt{0x0A..0x0F}}/{padding (NOPs)}}{
  \draw[thick, fill=amberbg] (0,\y-1.2) rectangle (2.6,\y);
  \node[align=center] at (1.3,\y-0.5) {\lbl};
  \node[align=center, font=\tiny, text=slate] at (1.3,\y-0.95) {\det};
}
\node[align=center, text=slate] at (1.3,0.0) {ROM 16 x 8 --- read-only,\\\texttt{pc} picks the row};
% data address space (right column)
\node[font=\small\bfseries] at (6.0,6.9) {data space};
\draw[thick] (3.6,0.4) rectangle (8.6,6.6);
\draw[thick, fill=goodbg] (3.6,5.0) rectangle (6.2,6.2);
\node[align=center] at (4.9,5.8) {regfile 8x8};
\node[align=center, font=\tiny, text=slate] at (4.9,5.35) {r0--r7 working registers\\(register-addressed, not memory)};
\draw[thick, fill=accentsoft] (6.4,5.0) rectangle (8.6,6.2);
\node[align=center] at (7.5,5.8) {RAM 16x8};
\node[align=center, font=\tiny, text=slate] at (7.5,5.35) {scratch data\\(BRAM-backed)};
\draw[thick, fill=amberbg] (3.6,3.4) rectangle (6.2,4.4);
\node[align=center] at (4.9,4.05) {flags reg};
\node[align=center, font=\tiny, text=slate] at (4.9,3.7) {Z N C V};
\draw[thick, fill=amberbg] (6.4,3.4) rectangle (8.6,4.4);
\node[align=center] at (7.5,4.05) {OUT reg};
\node[align=center, font=\tiny, text=slate] at (7.5,3.7) {to LEDs/UART};
\draw[thick, fill=graybg] (3.6,0.8) rectangle (8.6,2.9);
\node[align=center] at (6.1,2.2) {PC /4};
\node[align=center, font=\tiny, text=slate] at (6.1,1.5) {not in either space --- the sequencer\\lives beside memory, not inside it};
% note
\node[align=center, text=slate, text width=8.5cm] at (4.3,-0.7) {Two separate spaces: the program lives in ROM indexed by \texttt{pc}
(instruction fetch); everything a program can touch lives in the data space.
A von Neumann machine would share one space --- ours is Harvard-style.};
\end{tikzpicture}

\caption{The machine's map. Program space (ROM, pc-indexed) and data
space (registers, RAM, OUT) are separate --- a Harvard-style split.
The PC is not in either space: it is the sequencer that points
\emph{into} program space. Regenerate: \ttfamily
figures/ch15/fig\_memmap.tex.}
\label{fig:memmap15}
\end{figure}

\begin{figure}[htbp]
\centering
\begin{tikzpicture}[font=\scriptsize, >=latex]
% flops view (left)
\node[font=\small\bfseries] at (1.7,6.4) {as flops (what you wrote)};
\foreach \i in {0,...,3}{
  \pgfmathsetmacro{\y}{5.6-\i*1.1}
  \draw[thick, fill=graybg] (0.4,\y) rectangle (2.4,\y+0.8);
  \node at (1.4,\y+0.4) {8$\times$DFF};
}
\node[align=center, text=slate] at (1.4,0.9) {16 cells $\times$ 8 bits\\= up to 128 flops};
\draw[->, thick, slate] (3.2,3.0) -- (4.4,3.0);
% BRAM view (right)
\node[font=\small\bfseries] at (6.6,6.4) {as BRAM (what synthesis saw)};
\draw[thick, fill=goodbg] (4.6,1.8) rectangle (8.6,5.8);
\node[align=center] at (6.6,4.4) {\texttt{RAM16SDP4}};
\node[align=center, font=\tiny, text=slate] at (6.6,3.5) {simple dual-port\\1 write + 1 read port\\per cell};
\draw[->, thick] (3.8,4.6) -- (4.6,4.6) node[pos=-0.1, left]{\texttt{waddr,we}};
\draw[->, thick] (3.8,2.2) -- (4.6,2.2) node[pos=-0.1, left]{\texttt{raddr}};
\draw[->, thick, accent] (8.6,4.6) -- (9.4,4.6) node[right]{\texttt{wdata}};
\node[align=center, text=slate] at (6.6,0.9) {LUT fabric, not a flop file\\--- but still only 16x4 bits};
\node[align=center, text=slate, text width=9.5cm] at (4.8,-0.4) {Real
synth numbers: \texttt{regfile8} (two \emph{asynchronous} reads) $\to$
4$\times$\texttt{RAM16SDP4} distributed-RAM cells; \texttt{ram16}
(synchronous write \emph{and} read) $\to$ 2$\times$\texttt{RAM16SDP4}
plus 8 DFFs --- the output register is the extra flops. True block
BRAM needs the synchronous shape \emph{and} enough bits to be worth
the tile.};
\end{tikzpicture}

\caption{What synthesis actually did with our two memories ---
the real \texttt{synth\_gowin} numbers. The register file's
\emph{asynchronous} reads map onto distributed LUT-RAM cells
(\texttt{RAM16SDP4}); \texttt{ram16}'s all-synchronous shape uses half
as many cells plus eight DFFs --- the output register is exactly the
BRAM-style shape. True block BRAM is a bigger tile: it wants this
synchronous shape and thousands of bits to be worth spending.
Regenerate: \ttfamily figures/ch15/fig\_bram.tex.}
\label{fig:bram15}
\end{figure}

\section{Datapath versus control}\label{sec:control15}

\index{datapath}\index{control unit}
Here is the split that makes CPUs designable. The \emph{datapath} is
every register, memory, mux and ALU the data flows through --- the
``what can happen.'' The \emph{control unit} is the sequencer that
drives all the select lines and write-enables in the right order ---
the ``what happens now.'' This chapter builds the datapath and leaves
a hole where the control unit goes; ch16 fills the hole.

Our datapath needs eleven control bits:

\begin{center}
\begin{tabular}{ll}
\toprule
field & chooses\\
\midrule
\texttt{i\_alu\_op[2:0]} & which of the 8 ALU ops (ch14's contract)\\
\texttt{i\_wb\_sel[1:0]} & write-back source: ALU / imm byte / port A / port B\\
\texttt{i\_dst\_sel} & write destination: r0 (=A) or the operand register\\
\texttt{i\_we} & register-file write-enable\\
\texttt{i\_flags\_we} & latch Z\,N\,C\,V this step\\
\texttt{i\_out\_we} & latch the OUT register\\
\texttt{i\_pc\_mode[1:0]} & pc $+1$ / $+2$ / jump / hold\\
\bottomrule
\end{tabular}
\end{center}

Count them: $3+2+1+1+1+1+2=11$. That number is not trivia --- it is
the answer to ``how wide is the control word,'' the question every
control-unit design starts with.

\begin{exercise}[easy]
Which two control bits would you group into one field, and what
property of the ISA makes it safe? (Hint: what does \texttt{CMP}
share with \texttt{SUB}?)
\end{exercise}

\begin{figure}[htbp]
\centering
\begin{tikzpicture}[font=\scriptsize, >=latex, scale=0.88, transform shape]
% PC block
\draw[thick, fill=graybg] (0,5.6) rectangle (1.7,6.6);
\node[align=center] at (0.85,6.1) {PC /4\\(counter)};
% ROM
\draw[thick, fill=graybg] (3.0,5.6) rectangle (4.8,6.6);
\node[align=center] at (3.9,6.1) {ROM 16x8\\program};
\draw[->, thick] (1.7,6.1) -- (3.0,6.1) node[midway, above]{\texttt{pc}};
% instr fanout
\draw[->, thick] (4.8,6.1) -- (6.2,6.1) -- (6.2,5.0);
\draw[thick, fill=amberbg] (5.4,4.0) rectangle (7.0,5.0);
\node[align=center] at (6.2,4.5) {operand\\nibble /4};
% regfile
\draw[thick, fill=graybg] (2.0,1.6) rectangle (4.4,3.4);
\node[align=center] at (3.2,2.9) {regfile 8x8\\1W + 2R};
\node[align=center, font=\tiny, text=slate] at (3.2,2.1) {r0 = A};
\draw[->, thick] (6.2,4.0) -- (6.2,2.5) -- (4.4,2.5) node[pos=0.5, above]{\texttt{r[operand[2:0]]}};
% ALU
\draw[thick, fill=accentsoft] (5.6,0.6) rectangle (8.0,2.2);
\node[align=center] at (6.8,1.4) {\texttt{alu8}\\(ch14)};
\draw[->, thick] (4.4,2.0) -- (5.6,1.0) node[pos=0.4, below]{\texttt{A}};
\draw[->, thick] (4.4,1.8) -- (4.9,1.8) -- (4.9,0.2) -- (6.2,0.2) -- (6.2,0.6);
\node[below, font=\tiny] at (5.6,0.2) {\texttt{r[operand]}};
\draw[->, thick, accent] (7.2,2.2) -- (7.2,3.0);
\draw[thick, fill=goodbg] (6.3,3.0) rectangle (8.1,3.8);
\node[align=center] at (7.2,3.4) {flags reg\\Z N C V};
\draw[->, thick, accent] (8.1,3.4) -- (8.8,3.4) node[right]{\texttt{o\_z n c v}};
% WB mux
\draw[thick, fill=accentsoft] (1.4,-0.8) rectangle (3.2,-0.0);
\node[align=center] at (2.3,-0.4) {WB mux};
\draw[->, thick] (6.8,0.6) -- (6.8,-0.4) -- (3.2,-0.4);
\draw[->, thick] (2.3,0.0) -- (2.3,1.6) node[pos=0.5, left]{\texttt{wdata}};
\draw[->, thick] (4.8,6.4) -- (5.4,6.4) -- (5.4,7.0) -- (1.0,7.0) -- (1.0,-0.6) -- (1.4,-0.6) node[pos=0.15, above]{\texttt{imm byte}};
% OUT
\draw[thick, fill=goodbg] (9.4,0.6) rectangle (10.8,1.6);
\node[align=center] at (10.1,1.1) {OUT reg};
\draw[->, thick] (4.4,3.2) -- (9.8,3.2) -- (9.8,1.6) -- (10.1,1.6);
\draw[->, thick, accent] (10.8,1.1) -- (11.5,1.1) node[right]{LEDs/UART};
% control word input
\draw[thick, dashed, slate] (-0.4,-1.6) -- (11.5,-1.6) node[right, font=\tiny]{};
\node[align=left, text=slate] at (1.4,-2.0) {the 11-bit control word drives every dotted tap:\\
\texttt{alu\_op, wb\_sel, dst\_sel, we, flags\_we, out\_we, pc\_mode}};
\foreach \x/\y in {0.85/5.6, 3.9/5.6, 2.3/-0.8, 7.2/2.9, 10.1/0.6}{
  \draw[<-, dashed, slate] (\x,\y) -- (\x,\y-0.5);}
\draw[<-, dashed, slate] (6.8,0.6) -- (6.8,0.0);
\end{tikzpicture}

\caption{The datapath: the figure this book builds toward. PC feeds
the ROM; the instruction's operand nibble picks a register and the
immediate path; the ALU combines r0 (\texttt{A}) with the operand;
the write-back mux routes the result; flags and OUT latch on their
own enables. Every dotted tap is a bit of the 11-bit control word ---
which is all ch16's FSM has to drive. Regenerate: \ttfamily
figures/ch15/fig\_datapath.tex.}
\label{fig:datapath15}
\end{figure}

\section{The cycle: fetch, execute, write back}\label{sec:cycle15}

\index{fetch-execute cycle}\index{program counter}
A CPU repeats three moves forever: \emph{fetch} the instruction at
\texttt{pc}, \emph{execute} it through the datapath, \emph{write back}
whatever it produced (a register value, a flag, a pc change). In our
partial datapath the ``execute'' step is where the ALU, register file
and write-back mux do their work --- all in the same cycle, because
combinational logic settles before the edge.

The conventions we freeze here, so ch16 has no choices left to make:
the instruction byte is \texttt{\{op[7:4], operand[3:0]\}}; the byte
after it (\texttt{rom[pc+1]}) is an immediate when an op needs one;
\texttt{r0} doubles as the accumulator \texttt{A} that arithmetic ops
read and write; \texttt{MOV rd} stores A into a register; jumps take
their target from the operand nibble.

\begin{example}[Hand-tracing one full cycle]
Take step 5 of the demo program, \texttt{ADD r1} with A$=3$ and
r1$=5$:

\begin{center}
\begin{tabular}{lll}
\toprule
signal & driven & why\\
\midrule
\texttt{i\_alu\_op} & \texttt{ADD} & opcode 3 maps to ALU op 000\\
\texttt{i\_wb\_sel} & \texttt{ALU} & the result is the ALU's\\
\texttt{i\_dst\_sel} & \texttt{A} & accumulator ops write r0\\
\texttt{i\_we} & 1 & a real write-back\\
\texttt{i\_flags\_we} & 1 & arithmetic updates flags\\
\texttt{i\_pc\_mode} & \texttt{+1} & single-byte instruction\\
\bottomrule
\end{tabular}
\end{center}

On the clock edge: the ALU has already settled on $3+5=8$, the
write port latches it into r0, the flags register captures ZN CV
(0000), and pc becomes 7. Three registers changed on one edge ---
that \emph{is} the cycle.
\end{example}

\begin{figure}[htbp]
\centering
\begin{tikzpicture}[font=\scriptsize]
\node[anchor=north west, inner sep=0] at (0,0) {%
\begin{tabular}{llllll}
\toprule
step & \texttt{instr} & \texttt{wb\_sel} & \texttt{dst\_sel} & \texttt{pc\_mode} & writes\\
\midrule
1 & \texttt{LDI A,5} & \texttt{IMM} & \texttt{A} & \texttt{+2} & \texttt{r0$\leftarrow$0x05}\\
2 & \texttt{MOV r1}  & \texttt{PA}  & \texttt{reg} & \texttt{+1} & \texttt{r1$\leftarrow$A}\\
3 & \texttt{LDI A,3} & \texttt{IMM} & \texttt{A} & \texttt{+2} & \texttt{r0$\leftarrow$0x03}\\
4 & \texttt{MOV r2}  & \texttt{PA}  & \texttt{reg} & \texttt{+1} & \texttt{r2$\leftarrow$A}\\
5 & \texttt{ADD r1}  & \texttt{ALU} & \texttt{A} & \texttt{+1} & \texttt{r0$\leftarrow$8},\ flags\\
6 & \texttt{SUB r1}  & \texttt{ALU} & \texttt{A} & \texttt{+1} & \texttt{r0$\leftarrow$3},\ flags\\
7 & \texttt{OUT}     & ---      & ---      & \texttt{+1} & \texttt{out$\leftarrow$A}\\
8 & \texttt{HALT}    & ---      & ---      & \texttt{HOLD} & ---\\
\bottomrule
\end{tabular}};
\node[anchor=north west, text width=7.0cm, font=\scriptsize, text=slate]
  at (0,-5.3) {Every row is one clock edge and one 11-bit control word.
  \texttt{we}, \texttt{flags\_we} and \texttt{out\_we} are the
  write-enables the column hides: \texttt{we}=1 on rows 1--6,
  \texttt{flags\_we}=1 on rows 5--6 only, \texttt{out\_we}=1 on row 7.
  In ch16 an FSM emits this table from the opcode nibble; here the
  testbench plays it by hand --- the same sequence, two owners.};
\end{tikzpicture}

\caption{The whole demo program as a control-signal table --- one
control word per row, one clock edge per row. This table \emph{is}
what a control unit emits; the only difference between this chapter
and the next is who writes the rows. Regenerate: \ttfamily
figures/ch15/fig\_cycle.tex.}
\label{fig:cycle15}
\end{figure}

\begin{example}[The bench plays control unit]
\texttt{ch15\_datapath\_tb.v} drives that exact table, one control
word per step, and checks state after every edge. The real trace:

\begin{lstlisting}[caption={Real run of \ttfamily ch15\_datapath\_tb.v.
Watch \texttt{rb=xx} on the first MOV --- an unwritten register reads
\texttt{x}, which is the honest answer.}]
t=26000 pc= 2 instr=21 -> a=05 rb=xx flags z0 n0 c0 v0
t=36000 pc= 3 instr=10 -> a=05 rb=05 flags z0 n0 c0 v0
t=46000 pc= 5 instr=22 -> a=03 rb=xx flags z0 n0 c0 v0
t=56000 pc= 6 instr=31 -> a=03 rb=05 flags z0 n0 c0 v0
t=66000 ADD: a=08 z0 n0 c0 v0
PASS: datapath executed 8 control-word steps clean
\end{lstlisting}

Every line is a cycle the figure table predicted: pc marching through
ROM addresses, A being loaded and moved through registers, the ADD
producing 8 on the edge that also latches the flags. Nothing in this
trace is hypothetical --- it is the machine working.
\end{example}

\begin{example}[What it costs on the die]
Strategy: \texttt{synth\_gowin} the whole datapath and read the stat:

\begin{lstlisting}[caption={\ttfamily synth\_gowin stat for
\texttt{datapath} --- real output, trimmed.}]
=== datapath ===
      349 wires      610 wire bits       19 ports
      416 cells
       25   ALU          1   DFF         6   DFFR        16   DFFRE
        4   RAM16SDP4    62   LUT1       16   LUT2        30   LUT3
       81   LUT4       68   MUX2_LUT5   27   MUX2_LUT6  10   MUX2_LUT7
        3   MUX2_LUT8  13   IBUF        48   OBUF         4   $scopeinfo
\end{lstlisting}

Read it like the ch14 stat: \textbf{25 ALU} cells --- the two ALU
adders plus the pc incrementer. \textbf{23 flip-flops} (1 DFF + 6
DFFR + 16 DFFRE): the pc's 4 bits, the flags' 4 bits, the OUT
register's 8 --- the memories-with-enables, accounted for.
\textbf{4 RAM16SDP4} are the register file mapped onto distributed
RAM, as Figure~\ref{fig:bram15} explained. And \textbf{108 mux cells}
are the real price of the datapath pattern --- WB mux, dst select,
bus muxes: the wires are cheap, the choosers are the silicon.
\end{example}

\textbf{Honest marking.} The register-file, memory and datapath
benches and both synthesis stats were run and captured on this
machine. No board was attached; the datapath has no dedicated board
demo --- the full CPU gets one in ch16, where LEDs and UART have
something worth showing.

\section{Common mistakes}\label{sec:cm15}

\begin{tripwire}[Two talkers on one bus]
If two write sources can reach the same register input and both are
enabled in the same step, the hardware wires a fight. Rule: every
destination has \emph{one} write-back mux; every producer reaches it
through a select, never a parallel tie. If your design has two
drivers on one wire, there is no sim that can save you --- only the
architecture can.
\end{tripwire}
\begin{tripwire}[Expecting a real tri-state inside the FPGA]
\texttt{assign bus = oe ? val : 8'bz;} compiles; the synthesiser may
even map it. But there is no third state in the routing fabric ---
the tool silently turns it into pull-ups or muxes you did not design.
Draw the mux on purpose; never write \texttt{1'bz} inside a module's
internal wiring.
\end{tripwire}
\begin{tripwire}[Combinational reads where the silicon wants sync]
A register file reading asynchronously through a big array makes the
sim easy and the area ugly: async reads can't use block RAM at all.
Small files (ours is 64 bits) are fine as LUT-RAM; if you need a big
async-read array, the honest options are flops or a redesign --- not
hope.
\end{tripwire}
\begin{tripwire}[Forgetting a write-enable is a design choice]
\texttt{we=1} forever is how accidental designs happen: the file
updates every cycle whether the instruction asked or not. Every
state element in our datapath --- regfile, flags, OUT, pc --- has its
own enable, and \emph{defaults to not-writing}. Silence is the safe
state.
\end{tripwire}
\begin{tripwire}[Reading an unwritten register and believing it]
The demo's first \texttt{MOV r1} step reads \texttt{rb=xx} because r1
was never written. On a real machine, uninitialised registers are not
zero --- they are whatever they powered up as. Programs must write
before they read; testbenches should check that the \texttt{x} shows
up exactly where it should.
\end{tripwire}
\begin{tripwire}[Async inputs straight into a datapath]
A button or switch feeding a datapath signal can change \emph{during}
a clock edge --- metastability, the same lesson as ch10's
synchronisers. Any signal crossing into your clock domain passes
through a two-flop synchroniser first, or your clean datapath learns
analog electronics at the worst moment.
\end{tripwire}

\begin{exercise}[medium]
The demo program could not write to r4--r7 without changing encoding.
Say why (which instruction field is too narrow), and give the minimal
encoding change that frees the top four registers.
\end{exercise}
\begin{exercise}[medium]
\texttt{ram16} reads one edge after the address. If a program does
\texttt{LDI A; then uses the value}, how many wasted cycles does a
naive control unit burn per memory read, and what does the ch16 FSM
do with that cycle?
\end{exercise}
\begin{exercise}[hard]
Our pc counts bytes, so a two-byte LDI needs \texttt{pc+2}. Design
the alternative: pc counts \emph{instructions} (a small table maps
instruction index to byte address). What does it cost and what does
it fix?
\end{exercise}

\section{Summary}\label{sec:sum15}

\begin{center}
\begin{tabularx}{\linewidth}{@{}lX@{}}
\toprule
\textbf{Thing} & \textbf{The rule of thumb}\\
\midrule
register file & 1 clocked write port + N combinational read ports;
the write-enable is the ``remember'' bit\\
bus & one talker at a time: tri-state on breadboards, muxes in FPGA
fabric\\
program memory & instructions are just bytes; the pc is an address
counter reading them\\
BRAM inference & synchronous write \emph{and} read; async reads force
LUT-RAM or flops\\
datapath/control & datapath = what can happen; control = what happens
now --- the 11-bit control word is the seam\\
the cycle & fetch, execute, write-back --- one edge, three register
updates\\
\bottomrule
\end{tabularx}
\end{center}

\section*{Additional exercises}
\addcontentsline{toc}{section}{Additional exercises}

\begin{exercise}[easy]
Why does the register file's read side need no clock and no enable?
\end{exercise}
\begin{exercise}[easy]
How many bits wide is our datapath's control word, and which fields
does it split into?
\end{exercise}
\begin{exercise}[easy]
In the memory map, which space does the pc live in? Trick question
--- explain the answer.
\end{exercise}
\begin{exercise}[easy]
A register file is ``8$\times$8''. What do the two 8s count?
\end{exercise}
\begin{exercise}[medium]
Modify \texttt{regfile8} so r0 always reads as 0 (the riscv-zero
convention). Write the two-line change and say what it would break in
our accumulator-is-r0 convention.
\end{exercise}
\begin{exercise}[medium]
Design the control word for a \texttt{CMP rX} step (A$-rX$, flags
latched, nothing written). Which write-enables must be 0 and what
does that imply about the ALU's result?
\end{exercise}
\begin{exercise}[medium]
The tb drives control words on \texttt{negedge} and checks after
\texttt{posedge}. Explain the race this avoids, citing ch12.
\end{exercise}
\begin{exercise}[medium]
Widen the register file to 16$\times$8: which fields of the
instruction word break, and what is the minimal re-encoding that
keeps one-byte instructions?
\end{exercise}
\begin{exercise}[hard]
Write the truth table of \texttt{i\_pc\_mode}'s four encodings,
including what the FSM must guarantee when \texttt{HALT} holds
\texttt{pc} forever. What happens at power-on before the first HALT?
\end{exercise}
\begin{exercise}[hard]
Design the write-enable-trace bug we hinted at: a control table that
sets \texttt{we=1} during the OUT step. What does the regfile see,
and which tb check catches it first?
\end{exercise}
\begin{exercise}[hard]
Count the muxes in Figure~\ref{fig:datapath15} by hand, then compare
with the 108 \texttt{MUX2\_LUT*} cells in the stat report. Explain
the ratio --- what is the big mux the diagram hides?
\end{exercise}
\begin{exercise}[hard]
Our datapath reads the immediate byte through a second \texttt{rom16}
instance at \texttt{pc+1}. Give two reasons that's honest, and one
case where it lies.
\end{exercise}
\begin{exercise}[hard]
A teammate proposes dropping the flags register and reading ALU flag
outputs directly in the JZ step. Explain why that fails, referencing
when the flags are computed versus when they're sampled.
\end{exercise}
\begin{exercise}[hard]
Extend the datapath sketch with a stack pointer register and a
2-byte-deep call stack. Which existing fields handle
\texttt{CALL}/\texttt{RET} without new encoding space?
\end{exercise}
\begin{exercise}[hard]
\texttt{regfile8}'s async reads feed a combinational ALU feeding a
write-back mux --- a long combinational chain. Sketch the two-stage
version (operands registered) and name what it buys and what it costs
in control complexity.
\end{exercise}
\begin{exercise}[hard]
The stat report shows 48 OBUF on \texttt{datapath}. Reconcile that
number with the port list --- what are the 48, and why does a block
you'd never probe need outputs at all in a top-level stat?
\end{exercise}
\begin{exercise}[hard]
Write the \texttt{ram16} read-after-write hazard formally: if
\texttt{i\_we} and a read of the same address land on the same edge,
what does \texttt{o\_rdata} show and which BRAM policy does it
implement?
\end{exercise}
\begin{exercise}[hard]
The operand nibble is shared: register index \emph{and} jump target
\emph{and} part of the LDI immediate path. Prove it can't hold an
8-bit immediate in one byte, and state the design rule that keeps
encoding pressure sane.
\end{exercise}

\section*{Hints and answers}
\addcontentsline{toc}{section}{Hints and answers}
\begin{description}
\item[Control-bit grouping (easy, inline)] \texttt{dst\_sel} folds
into \texttt{wb\_sel} if arithmetic ops always write A --- ours do;
a CMP's ``no write'' already lives in \texttt{we}.
\item[Why reads are free (easy)] Read ports are mux outputs: address
in, value out, nothing stored --- only writes need an edge and an
enable.
\item[Control word (easy)] 11 bits: alu\_op3 + wb\_sel2 + dst\_sel1 +
we + flags\_we + out\_we + pc\_mode2.
\item[PC's space (easy)] Neither --- the pc is an index
\emph{into} program space, part of the sequencer, not of memory.
\item[8$\times$8 (easy)] 8 registers deep, 8 bits wide each.
\item[r0-hardwire (medium)] \texttt{assign o\_ra = (addr==0) ? 0 :
mem[addr]} --- it would break \texttt{A}=r0, because A could never
hold a value.
\item[CMP word (medium)] alu\_op=CMP, wb\_sel=ALU, we=\textbf{0},
flags\_we=1, pc=+1 --- the result computes and dies on the mux;
flags are the verdict.
\item[negedge drive (medium)] driving stimulus on the opposite edge
avoids the tb/DUT race of ch12: DUT samples old values at posedge,
tb's new values apply after.
\item[16$\times$8 (medium)] operand[3:0] can already name 0--15 ---
widen \texttt{i\_raddr}/\texttt{i\_waddr} to 4 bits; the encoding
survives, only registers 8+ become usable.
\item[pc\_mode truth (hard)] HOLD must be reachable \emph{after} the
halt instruction lands, or pc wraps to 0 and re-executes --- the
power-on path is reset$\to$fetch, the same first step every time.
\item[OUT we bug (hard)] the regfile latches wb\_data wherever dst
points --- usually A gets overwritten with alu\_y; the post-OUT A
check fires first.
\item[Mux census (hard)] the diagram shows ~5 muxes (WB, dst, B-src,
pc-next, flag-sel); 108 LUT muxes is the \emph{bit-level} count ---
each of our 8-bit muxes is eight 1-bit muxes deep.
\item[Second rom16 (hard)] honest because ROM reads are pure and the
immediate \emph{is} the next byte; it lies if an op's operand is
meant as an address, not data --- the second instance can't fetch
operands from data space.
\item[No flags reg (hard)] ALU flags vanish when the inputs change;
JZ reads them a step after the arithmetic settled --- unregistered
flags would reflect the wrong operation.
\item[Stack (hard)] \texttt{CALL} = pc\_mode jump + write pc+1 to
stack; \texttt{RET} = jump to stack top --- fits in pc\_mode plus one
stack-pointer we.
\item[Two-stage (hard)] registers operands on an edge, ALU works next
cycle: halves the combinational depth, but every op takes 2 cycles
and the control table doubles.
\item[48 OBUF (hard)] pc4 + instr8 + imm8 + a8 + rb8 + out8 + flags4
= 48: in a top-level stat, debug outputs cost real pins --- strip
them for the CPU build or let hierarchy absorb them.
\item[RAW hazard (hard)] \texttt{o\_rdata} shows the \emph{old} value
--- read-before-write, the BRAM default; write-first policies exist
but ours didn't ask for one.
\item[Nibble budget (hard)] operand[3:0] has 4 bits; imm8 needs 8 ---
the LDI immediate must come from the next byte, which is exactly
why pc has a +2 mode: fetch two bytes per instruction that needs it.
\end{description}

\section*{Further reading}
\addcontentsline{toc}{section}{Further reading}

\begin{itemize}
\item Harris \& Harris, \emph{Digital Design and Computer
Architecture}, 2nd ed., ch.~7 (microarchitecture).\newline
\url{pages.hmc.edu/harris/ddca/} --- the single-cycle datapath our
figures follow; the control/datapath split is theirs.
Intermediate. Verified 200, 2026-10-11.
\item Nisan \& Schocken, \emph{The Elements of Computing Systems},
projects 3--5 (RAM, registers, the Hack computer).\newline
\url{nand2tetris.org/project03} and \url{nand2tetris.org/project05}
--- the ladder from DFF to a working machine; their control-word
table inspired Figure~\ref{fig:cycle15}. Intermediate. Verified 200,
2026-10-11.
\item \emph{HDLBits} --- \url{hdlbits.01xz.net} ---
Circuits/Sequential ``Latches and Flip-Flops'' through ``Shift
Registers'' drills the register idioms this chapter builds on.
Beginner. Verified 200, 2026-10-11.
\item Lushay Labs, \emph{Our First CPU} (Tang Nano series).
\newline\url{learn.lushaylabs.com/tang-nano-series} --- a register
file and datapath on this exact board, with working code.
Beginner--intermediate. Verified 200, 2026-10-11.
\item ChipVerify, \emph{Verilog tutorial} --- memory arrays, FSM
glue and the register-transfer idioms used here.\newline
\url{chipverify.com/tutorials/verilog} --- Beginner. Verified 200,
2026-10-11.
\item Wikipedia, \emph{Datapath} and \emph{Register file} ---
terminology and the wider taxonomy (multi-issue, banks, forwarding)
this small machine points toward. Reference. Verified 200,
2026-10-11.\newline\url{en.wikipedia.org/wiki/Datapath}\newline
\url{en.wikipedia.org/wiki/Register_file}
\end{itemize}

\begin{rulebox}[{Watch instead}]
\begin{enumerate}
\item \emph{Bus architecture and how register transfers work ---
8-bit register, Part 1} --- Ben Eater, 9:35,
\url{https://www.youtube.com/watch?v=QzWW-CBugZo} --- the bus/tri-state
section taught on real breadboard wiring; watch the register transfer
happen on the bus before reading our mux version. Matches
\S\ref{sec:bus15}--\S\ref{sec:control15}. Opened and confirmed
2026-10-11.
\item \emph{Registers and RAM: Crash Course Computer Science \#6}
--- CrashCourse, 12:16,
\url{https://www.youtube.com/watch?v=fpnE6UAfbtU} --- from latch to
register file to addressable memory, the same ladder in 12 minutes.
Matches \S\ref{sec:regfile15} and \S\ref{sec:memory15}. Opened and
confirmed 2026-10-11.
\end{enumerate}
\end{rulebox}

\chapternotes{The datapath/control split and the single-cycle
structure follow Harris~\&~Harris ch.~7~\cite{ddca}; the
register-file interface and control-word table follow the
nand2tetris projects~3--5 conventions. Memory-inference numbers are
real \texttt{yosys} 0.69 \texttt{synth\_gowin} runs
\cite{apicula}; the Harvard split (separate program/data spaces) is
the book's own encoding decision, seeded here and completed by the
ch16 ISA table.}
